% =====================================================================
%  exemplo.tex -- Documento de exemplo do sistema visual PousadaGuide
%  Compile com:  latexmk -xelatex exemplo.tex     (ou: tectonic exemplo.tex)
% =====================================================================
\documentclass{pousadaguide}
% \documentclass[print]{pousadaguide}   % versão para impressão (fundo branco)

% ---- Metadados do guia ---------------------------------------------
\guianumero{00}
\guiatitulo{Guia de Exemplo}
\guiasubtitulo{Um tour por todos os componentes do sistema visual dos PousadaGuides.}
\guiaslug{exemplo}
\guiacategoria{sistema visual}
\guiaversao{v1.0}
\guiaautor{Bakalhau}
\guiadata{novembro de 2026}
\guialicenca{GNU AGPLv3}
\guialinks{%
  \textcolor{pgGlow}{\faDiscord}\ \ discord.gg/pousada\\[1mm]
  \textcolor{pgGlow}{\faGithub}\ \ github.com/pousada\\[1mm]
  \textcolor{pgGlow}{\faRss}\ \ pousada.substack.com}

\begin{document}

\capa
\tableofcontents

% =====================================================================
\comentario{antes de tudo, entenda o terreno}
\chapter{Primeiros passos}

\begin{missao}
\begin{itemize}
  \item Entender como um PousadaGuide é estruturado.
  \item Conhecer as caixas de destaque e quando usar cada uma.
  \item Rodar seu primeiro comando no terminal com segurança.
\end{itemize}
\end{missao}

Este é o texto corrido do guia, usando \textbf{IBM Plex Sans} para leitura
confortável em tela. Títulos, rótulos e código usam \textbf{IBM Plex Mono}, que
dá o ar ``terminal'' da identidade. Você pode usar \emph{ênfase}, links como
\pglink{https://fmhy.net}, links com texto como \pglink[o repositório]{https://github.com},
código em linha como \cmd{sudo pacman -Syu} e atalhos como
\kbd{Ctrl}+\kbd{Shift}+\kbd{Del}.

\section{Caixas de destaque}

Use as caixas com moderação: elas chamam atenção justamente por serem raras.

\begin{nota}
Informação neutra, contexto ou curiosidade que complementa o texto.
\end{nota}

\begin{dica}[atalho útil]
Boas práticas e truques que economizam tempo. O texto entre colchetes vira o
título da caixa.
\end{dica}

\begin{atencao}
Algo que pode dar errado se o leitor não prestar atenção, mas sem risco grave.
\end{atencao}

\begin{perigo}[nunca faça isso]
Risco real de segurança, perda de dados ou exposição. Use só quando for sério.
\end{perigo}

\begin{tldr}
Resumo de uma seção longa em duas ou três linhas, para quem está com pressa.
\end{tldr}

\begin{pousadinha}
Oi! Eu apareço quando o guia quer falar mais de perto com você, explicar algo
de um jeito mais leve ou dar aquela dica de quem já quebrou a cara antes.
\end{pousadinha}

\section{Terminal e código}

\subsection{Comandos de shell}

O \texttt{\$} fica destacado automaticamente e comentários com \texttt{\#}
ficam esmaecidos:

\begin{terminal}[bash -- ~/downloads]
# atualiza o sistema (Arch / Manjaro)
$ sudo pacman -Syu

# confere o hash de um arquivo baixado
$ sha256sum ubuntu-24.04-desktop-amd64.iso
a1b2c3d4e5f6...  ubuntu-24.04-desktop-amd64.iso
\end{terminal}

\subsection{Arquivos de código}

\begin{codigo}[language=JavaScript]{exemplo.user.js}
// ==UserScript==
// @name         Olá Pousada
// @match        https://*/*
// ==/UserScript==
(function () {
  'use strict';
  const msg = "Bem-vindo à Pousada!"; // comentário com acentuação
  console.log(msg);
})();
\end{codigo}

\subsubsection{Nível mais profundo de título}

Subseções usam \texttt{>>} e o terceiro nível usa \texttt{//}, como comentários
de código.

% =====================================================================
\comentario{o arsenal recomendado pela comunidade}
\chapter{Ferramentas}

\section{Fichas de ferramenta}

\begin{ferramenta}[recomendado]{Zen Browser}{https://zen-browser.app}{open source, firefox, win, linux, mac}
Navegador baseado no Firefox, focado em privacidade e com uma interface limpa
com abas verticais. \estrelas{5}
\end{ferramenta}

\begin{ferramenta}{qBittorrent}{https://www.qbittorrent.org}{open source, torrent, win, linux, mac}
Cliente de torrent livre, sem anúncios, com busca integrada e suporte a filtros
de IP. \estrelas{4}
\end{ferramenta}

\section{Comparativos}

\begin{tabela}{>{\bfseries}l X c c c}
\cab{Cliente} & \cab{Observação} & \cab{Livre} & \cab{Anúncios} & \cab{Leve} \\
qBittorrent  & Melhor equilíbrio geral          & \ok & \ok   & \meio \\
Transmission & Minimalista, ótimo no Linux/macOS & \ok & \ok   & \ok   \\
Deluge       & Extensível via plugins           & \ok & \ok   & \ok   \\
µTorrent     & Evite: anúncios e histórico ruim & \nope & \nope & \meio \\
\end{tabela}

Etiquetas avulsas: \pgtag{open source} \pgtag[pgAmber]{freemium}
\pgtag[pgAlert]{evite} \pgtag[pgBrass]{beta}

\section{Passo a passo}

\begin{passo}{Baixe o instalador oficial}
Sempre pelo site oficial do projeto ou pelo gerenciador de pacotes do seu sistema.
\end{passo}

\begin{passo}{Verifique a integridade}
Compare o hash \cmd{SHA-256} com o publicado na página de download.
\end{passo}

\begin{passo}{Instale e configure}
Siga as configurações recomendadas da próxima seção.
\end{passo}

\section{Listas}

\begin{itemize}
  \item Itens de lista usam o prompt \texttt{>} como marcador.
  \begin{itemize}
    \item E o segundo nível usa um traço.
  \end{itemize}
  \item Simples e legível.
\end{itemize}

\begin{enumerate}
  \item Listas numeradas usam dois dígitos.
  \item Como contadores de terminal.
\end{enumerate}

\begin{checklist}
  \feito uBlock Origin instalado
  \feito DNS com bloqueio configurado
  \item VPN configurada
  \opcional Debrid contratado
\end{checklist}

\divisor

\begin{description}
  \item[Debrid] Serviço que baixa o arquivo por você em servidores rápidos.
  \item[Seed] Quem continua compartilhando o arquivo depois de baixar.
\end{description}

\contracapa
\end{document}
